\section{Анализ предметной области}
\subsection{Эволюция систем производственного планирования}
\subsubsection{От ручных графиков к автоматизированным системам}
Задача упорядочения работ и распределения ресурсов возникла одновременно с появлением мануфактур. Первые методы планирования опирались на бумажные диаграммы Ганта и интуицию мастеров. Диаграмма, предложенная Генри Гантом в 1910‑х годах, долгое время оставалась основным инструментом визуализации производственных процессов. Её главное достоинство — наглядность, а недостаток — полное отсутствие оптимизации: любое изменение загрузки или поломка станка требовали ручного пересчёта всего графика.

С ростом сложности изделий и числа операций ручное планирование перестало справляться. В 1960–1970‑х годах появились первые системы класса MRP (Material Requirements Planning), которые автоматизировали расчёт потребностей в материалах и комплектующих. Они работали по жёсткому алгоритму, не учитывая ограничения по мощностям. Это породило следующее поколение — MRP II (Manufacturing Resource Planning), где к материалам добавились производственные мощности и финансы. Однако MRP II всё ещё предполагал идеальные условия и не мог динамически реагировать на изменения в цехе.

\subsubsection{Появление ERP-систем и их роль в планировании}
В 1990‑х годах концепция MRP II эволюционировала в ERP (Enterprise Resource Planning) — системы управления всеми ресурсами предприятия. ERP объединили финансы, логистику, производство, кадры и продажи в единой базе данных. Наиболее известные представители — SAP ERP, Oracle E-Business Suite, Microsoft Dynamics, а также российские 1С:ERP и Галактика ERP.

Производственный модуль ERP позволяет формировать планы производства на основе заказов клиентов и прогнозов спроса, отслеживать выполнение операций и вести учёт затрат. Однако стандартные ERP-системы решают задачу планирования упрощённо: они строят план «назад» от даты отгрузки, предполагая бесконечную мощность ресурсов. Это приводит к нереалистичным расписаниям в цехах с плотной загрузкой и переналадками. Кроме того, ERP не учитывает фактическое состояние оборудования в реальном времени, что требует дополнительного уровня управления.

\subsubsection{Специализированные APS-системы}
Ограниченность ERP-модулей вызвала появление Advanced Planning \& Scheduling (APS) — продвинутых систем планирования и диспетчеризации. APS работают как надстройка над ERP или MES, получая от них заказы и данные о ресурсах, а возвращая выполнимое расписание с точностью до минуты. Ключевые возможности APS:
\begin{itemize}
    \item одновременный учёт материалов и мощностей;
    \item оптимизация по заданным критериям (минимальное опоздание, максимальная загрузка);
    \item поддержка альтернативных маршрутов и ресурсов;
    \item возможность перепланирования «на лету» при возникновении отклонений.
\end{itemize}
Промышленные APS-решения (например, Siemens Opcenter APS, Asprova, Preactor) используют точные методы оптимизации, эвристики или их комбинацию. Их общий недостаток — высокая стоимость лицензий и длительное время расчёта для крупных цехов.

\subsubsection{Индустрия 4.0 и искусственный интеллект в планировании}
Четвёртая промышленная революция принесла цифровые двойники, интернет вещей и массовое применение машинного обучения. Нейросетевые диспетчеры способны обучаться на исторических данных или синтетических примерах и выдавать решение за миллисекунды. Это открывает путь к по-настоящему адаптивному управлению производством, где расписание корректируется в реальном времени при поступлении срочного заказа или поломке станка. Интеллектуальные APS-системы способны не только оптимизировать текущее расписание, но и прогнозировать узкие места, рекомендовать профилактическое обслуживание и моделировать различные сценарии загрузки.

\subsection{Основы производственного планирования и диспетчеризации}

\subsubsection{Задача составления расписаний}
В теории расписаний задача Job Shop Scheduling (JSS) формулируется как распределение  работ по  машинам с соблюдением технологической последовательности операций. Классический JSS является NP-трудным: число возможных расписаний растёт факториально, и точное решение для десятков операций требует неприемлемого времени. Для небольших задач (до 10–15 операций) применяются точные методы, такие как ветвей и границ или целочисленное линейное программирование (MILP). Однако с ростом размерности их вычислительная сложность делает их непригодными для оперативного планирования.

\subsubsection{Гибкое производственное расписание (Flexible Job Shop)}
В реальном производстве операция может выполняться не на одном, а на нескольких альтернативных станках, причём с разной длительностью. Это задача Flexible Job Shop Scheduling (FJSS). Она добавляет ещё один уровень сложности: помимо определения порядка запуска, необходимо выбрать подходящий ресурс. Именно этот класс задач является целевым для разрабатываемой системы. FJSS формулируется как две взаимосвязанные подзадачи: маршрутизация (выбор ресурса) и секвенирование (определение порядка операций на каждом ресурсе). Обе подзадачи NP-трудны, что делает FJSS одной из самых сложных задач комбинаторной оптимизации.

\subsubsection{Ключевые показатели эффективности расписаний}
Качество расписания оценивается набором KPI:
\begin{itemize}
    \item Makespan — общее время выполнения всех работ (от старта первой до завершения последней операции);
    \item Total tardiness — суммарное опоздание по всем заказам;
    \item Mean tardiness — среднее опоздание заказа;
    \item Percent of tardy jobs — доля опоздавших заказов;
    \item Resource utilization — коэффициент загрузки оборудования;
    \item Balance (CV of load) — равномерность распределения нагрузки.
\end{itemize}
Различные методы планирования могут оптимизировать один или несколько показателей. На практике часто ищут компромисс между минимальным опозданием и равномерной загрузкой.

\subsubsection{Неопределённость и человеческий фактор}
В реальном цехе нормативные длительности операций редко совпадают с фактическими, станки выходят из строя, поступают срочные заказы, отсутствуют операторы. Это делает статическое расписание, построенное утром, неактуальным уже к полудню. Отсюда потребность в оперативном перепланировании — способности системы быстро перестроить график с учётом текущей обстановки. Традиционные подходы, основанные на периодическом перезапуске MILP, не успевают реагировать на события в реальном времени. Нейросетевые методы, наоборот, позволяют мгновенно пересчитывать назначения, что делает их идеальными кандидатами для динамических сред.

\subsubsection{Классификация методов решения FJSS}
Все методы решения FJSS можно разделить на точные, эвристические и метаэвристические. Точные методы (MILP, динамическое программирование) гарантируют оптимальность, но экспоненциально растут по времени. Эвристики (диспетчерские правила) работают быстро, но не гарантируют качества. Метаэвристики (генетические алгоритмы, муравьиные колонии, имитация отжига) ищут компромисс. Нейросетевые методы занимают особое место: они обучаются на примерах (возможно, полученных от точных методов или эвристик) и затем работают с постоянной скоростью, близкой к эвристикам, но с качеством, приближающимся к метаэвристикам.

\subsection{Системы управления производством: MES и ERP}

\subsubsection{ERP-системы}
ERP решают задачи верхнего уровня: приём заказов, расчёт потребностей в сырье, планирование производства на дни и недели. Они оперируют понятиями «заказ», «партия», «маршрутная карта», но не знают деталей цехового исполнения. Результат ERP-планирования — укрупнённый план, который должен быть детализирован и адаптирован к фактической ситуации в цехе. Исторически ERP развивались из бухгалтерских систем, что наложило отпечаток на их архитектуру: они прекрасно считают деньги, но плохо управляют временем с точностью до минуты.

\subsubsection{MES-системы}
Manufacturing Execution System (MES) действует на цеховом уровне. Она собирает данные о состоянии оборудования, фиксирует фактическое выполнение операций, управляет очередями, формирует задания операторам. MES является связующим звеном между верхнеуровневым планом ERP и реальным производством. Однако стандартные MES не содержат мощных оптимизационных алгоритмов и часто полагаются на простые правила диспетчеризации (FIFO, EDD, SPT). Их задача — исполнение плана, а не его оптимизация.

\subsubsection{Интеграция ERP/MES/APS}
Эффективная архитектура предполагает, что ERP передаёт заказы в APS, APS строит детальное расписание и отдаёт его в MES на исполнение. MES собирает фактические данные и сообщает об отклонениях, на основе которых APS перестраивает план. Такая трёхзвенная модель обеспечивает сквозное управление производством от заказа клиента до отгрузки. На практике интеграция осложняется разными форматами данных, разными поставщиками систем и необходимостью синхронизации мастер-данных.

\subsubsection{Примеры программных продуктов}
В таблице \ref{software_compare:table} приведён обзор существующих решений для планирования производства, не использующих нейронные сети.

\renewcommand{\arraystretch}{0.8} % уменьшение расстояний до сетки таблицы
\begin{xltabular}{\textwidth}{|p{2.5cm}|X|>{\setlength{\baselineskip}{0.7\baselineskip}}p{2.85cm}|>{\setlength{\baselineskip}{0.7\baselineskip}}p{4.85cm}|}
\caption{Сравнение промышленных APS-решений \label{software_compare:table}}\\
\hline \centrow \setlength{\baselineskip}{0.7\baselineskip} Продукт& \centrow \setlength{\baselineskip}{0.7\baselineskip} Метод оптимизации & \centrow Интеграция & \centrow Особенности \\
\hline \centrow 1 & \centrow 2 & \centrow 3 & \centrow 4\\ 
\hline
\endfirsthead
\continuecaption{Продолжение таблицы \ref{software_compare:table}}\\
\hline \centrow 1 & \centrow 2 & \centrow 3 & \centrow 4\\ 
\hline
\finishhead
Siemens Opcenter APS & Генетический алгоритм + эвристики & ERP/MES Siemens & Высокая стоимость, сложное внедрение \\ 
\hline Asprova APS & Эвристики с настраиваемыми правилами & Различные ERP & Поддержка «drag-and-drop» ручной коррекции \\ 
\hline Preactor (Siemens) & Дискретно-событийное моделирование & Собственное API & Гибкое моделирование, но медленное для больших задач \\ 
\hline 1С:ERP + MES & Правила приоритетов, MRP II & Внутренняя экосистема 1С & Доступность для российских предприятий
\end{xltabular}
\renewcommand{\arraystretch}{1.0} % восстановление сетки

\subsubsection{Сравнение централизованной и децентрализованной диспетчеризации}
Традиционные APS централизованно решают задачу для всего цеха, что требует полной информации о всех заказах и ресурсах. В противоположность этому, децентрализованные (мультиагентные) подходы делегируют принятие решений самим ресурсам или операциям. Агенты «договариваются» между собой, что повышает живучесть системы, но может приводить к субоптимальным решениям. Нейросетевой диспетчер, реализованный в данной работе, занимает промежуточную позицию: решение принимается централизованно, но архитектура (Pointer Network) позволяет естественным образом обрабатывать переменное число агентов-ресурсов.

\subsection{Нейросетевые методы в планировании и диспетчеризации}

\subsubsection{Обзор подходов}
Современные исследования предлагают два основных направления применения нейронных сетей в планировании:
\begin{itemize}
    \item Обучение с учителем (поведенческое клонирование): модель учится копировать решения экспертной системы или MILP-солвера. После обучения она способна мгновенно повторять «хорошие» назначения.
    \item Обучение с подкреплением (RL): агент взаимодействует со средой (симулятором цеха), получая награду за своевременное выполнение заказов. Со временем он вырабатывает стратегию, максимизирующую суммарную награду.
\end{itemize}
Промежуточный вариант — самоконтролируемое обучение (self‑labeling), когда модель сначала обучается на размеченных данных, а затем сама генерирует дополнительные примеры, в которых она уверена, и доучивается на них. Такой подход позволяет улучшить обобщающую способность модели без привлечения дополнительных размеченных данных.

\subsubsection{Архитектуры нейронных сетей}
Для задачи выбора ресурса в FJSS применяются:
\begin{itemize}
    \item Многослойный перцептрон (MLP) с фиксированным числом входов — подходит, когда число ресурсов постоянно;
    \item Pointer Network и механизмы внимания (Attention) — позволяют обрабатывать переменное число ресурсов, указывая на наиболее подходящий;
    \item Графовые нейронные сети (GNN) — кодируют структуру цеха в виде графа.
\end{itemize}
В разрабатываемой системе используется комбинация MLP для фиксированной конфигурации и Pointer Network как масштабируемое решение. MLP хорошо работает, когда число ресурсов известно заранее и не меняется. Pointer Network (или аналогичные attention-модели) могут принимать на вход последовательность ресурсов переменной длины и выбирать один из них, что делает архитектуру инвариантной к числу станков.

\subsubsection{Преимущества нейросетевых диспетчеров}
\begin{itemize}
    \item Скорость: расписание строится за доли секунды.
    \item Адаптивность: модель может быть дообучена на новых данных при изменении конфигурации цеха.
    \item Компактность: обученная модель весит мегабайты и не требует мощного сервера.
\end{itemize}
Ограничения: модель не гарантирует оптимальности; качество зависит от репрезентативности обучающей выборки. Кроме того, модель может «забывать» старые паттерны при дообучении на новых данных (проблема катастрофического забывания). Для борьбы с этим могут применяться методы инкрементального обучения или rehearsal-стратегии.

\subsubsection{Примеры внедрения}
В научной литературе описан ряд успешных экспериментов, где нейросетевой диспетчер сокращал среднее опоздание на 15–25\% по сравнению с классическими эвристиками, сохраняя стабильность при изменении числа заказов и ресурсов. Подобные результаты подтверждают перспективность подхода, реализованного в данной работе. Например, в статье \textit{Zhang et al. (2020)} показано, что Pointer Network, обученный на решениях генетического алгоритма, достигает 98\% качества эталона при тысячекратном ускорении. В работе \textit{Kim et al. (2021)} RL-агент на основе PPO сократил makespan на 12\% по сравнению с эвристикой EDD.
                     %                                                                                       ^^^^^^^^^^^^^^^ -           ТУТ НАДО ССЫЛКУ
\subsubsection{Сравнение методов обучения}
В таблице \ref{learning_compare:table} приведено сравнение трёх парадигм обучения, реализованных в системе.

\renewcommand{\arraystretch}{0.8} % уменьшение расстояний до сетки таблицы
\begin{xltabular}{\textwidth}{|X|p{2.5cm}|>{\setlength{\baselineskip}{0.7\baselineskip}}p{2.85cm}|>{\setlength{\baselineskip}{0.7\baselineskip}}p{2.85cm}|}
\caption{Сравнение методов обучения нейросетевых диспетчеров \label{learning_compare:table}}\\
\hline \centrow \setlength{\baselineskip}{0.7\baselineskip} Характеристика& \centrow \setlength{\baselineskip}{0.7\baselineskip} BC (PointerNet) & \centrow SLIM (Self‑Labeling) & \centrow PPO (RL) \\
\hline \centrow 1 & \centrow 2 & \centrow 3 & \centrow 4\\ 
\hline
\endfirsthead
\continuecaption{Продолжение таблицы \ref{software_compare:table}}\\
\hline \centrow 1 & \centrow 2 & \centrow 3 & \centrow 4\\ 
\hline
\finishhead
Требует размеченные данные & Да & Да (начально) & Нет \\ 
\hline Способность к обобщению & Средняя & Высокая & Высокая \\ 
\hline Время обучения & Низкое & Среднее & Высокое \\ 
\hline Качество на тестовом сценарии & Хорошее & Наилучшее & Наихудшее \\ 
\hline Устойчивость к шуму в данных & Низкая & Средняя & Высокая
\end{xltabular}
\renewcommand{\arraystretch}{1.0} % восстановление сетки

\subsection{Целесообразность внедрения интеллектуальной APS-системы}

\subsubsection{Критерии эффективности}
Решение о внедрении интеллектуального планировщика должно опираться на количественные показатели. Основные критерии:
\begin{itemize}
    \item сокращение времени построения расписания;
    \item уменьшение среднего опоздания заказов;
    \item рост коэффициента загрузки оборудования;
    \item снижение доли опаздывающих заказов;
    \item равномерность распределения нагрузки.
\end{itemize}
В таблице \ref{before_after:table} приведены данные для цеха из 5 станков и 120 заказов, демонстрирующие эффект от внедрения.

\renewcommand{\arraystretch}{0.8} % уменьшение расстояний до сетки таблицы
\begin{xltabular}{\textwidth}{|X|p{2.5cm}|>{\setlength{\baselineskip}{0.7\baselineskip}}p{2.85cm}|}
\caption{Сравнение «до» и «после» внедрения нейросетевого планировщика \label{before_after:table}}\\
\hline \centrow \setlength{\baselineskip}{0.7\baselineskip}Показатель& \centrow \setlength{\baselineskip}{0.7\baselineskip}Эвристика & \centrow SLIM-модель \\
\hline \centrow 1 & \centrow 2 & \centrow 3\\ 
\hline
\endfirsthead
\continuecaption{Продолжение таблицы \ref{software_compare:table}}\\
\hline \centrow 1 & \centrow 2 & \centrow 3\\ 
\hline
\finishhead
Среднее опоздание, ч & 9.0 & 8.1 \\ 
\hline Максимальное опоздание, ч & 31.2 & 23.0 \\ 
\hline Доля опаздывающих заказов, \% & 85 & 85 \\ 
\hline Makespan, ч & 43.2 & 35.0 \\ 
\hline Коэф. вариации загрузки & 0.314 & 0.171
\end{xltabular}
\renewcommand{\arraystretch}{1.0} % восстановление сетки

Как видно из таблицы, нейросетевой диспетчер улучшает makespan на 19\% и делает загрузку в 1.8 раза равномернее. Экономический эффект достигается за счёт более полного использования дорогостоящего оборудования и сокращения штрафов за срыв сроков.

\subsubsection{Технические и организационные аспекты внедрения}
Для успешного внедрения требуется:
\begin{itemize}
    \item наличие исторических данных или экспертных правил для обучения модели;
    \item интеграция с существующей ERP/MES (при необходимости);
    \item обучение персонала работе с новым интерфейсом;
    \item постепенное расширение функциональности — от параллельной работы с существующим планировщиком до полного замещения.
\end{itemize}
Разработанная система удовлетворяет этим требованиям: она автономна, не требует внешних баз данных и может быть развёрнута на обычном рабочем компьютере.

\subsubsection{Перспективы развития}
В дальнейшем система может быть расширена:
\begin{itemize}
    \item поддержкой динамических событий в реальном времени (поломка, срочный заказ);
    \item интеграцией с IoT-датчиками для автоматического определения состояния станков;
    \item переходом на полностью автономное управление цехом с минимальным участием человека.
\end{itemize}
Таким образом, разработанный программный продукт не только решает текущие задачи планирования, но и закладывает фундамент для создания «умного» производства следующего поколения.

\subsection{Анализ производственных данных и потоков работ}
\subsubsection{Жизненный цикл производственного заказа}
Производственный заказ проходит несколько этапов: регистрация в ERP, разузлование (определение состава операций и материалов), планирование (APS), исполнение (MES) и закрытие. На каждом этапе могут возникать задержки и отклонения. Задача интеллектуальной системы — минимизировать задержки именно на этапе планирования и диспетчеризации, обеспечивая максимально гладкий поток работ.

\subsubsection{Типовая структура производственного цеха}
Цех состоит из набора рабочих центров (станков), каждый из которых может выполнять один или несколько типов операций. Между станками могут существовать различия в скорости обработки, точности, надёжности. Некоторые станки могут быть взаимозаменяемыми, другие — уникальными. Все эти особенности должны быть отражены в модели данных системы.

\subsubsection{Потоки данных между системами}
На рисунке \ref{potoki:image} схематично показаны потоки данных между ERP, APS, MES и уровнем оборудования. 

\begin{figure}[H]
\center{\includegraphics[width=1\linewidth]{potoki}}
\caption{Cхема потоков данных основных систем предприятия}
\label{potoki:image}
\end{figure}

Планирование начинается с получения заказов из ERP. APS обогащает их технологическими картами, учитывает доступность ресурсов и строит расписание. MES получает расписание и доводит его до операторов и станков, фиксируя факт выполнения. Отклонения возвращаются в APS для перепланирования.

\subsubsection{Роль человека в автоматизированном планировании}
Несмотря на высокую степень автоматизации, человек остаётся ключевым звеном: он может корректировать приоритеты, переназначать операции вручную, подтверждать или отклонять предложенные планы. Поэтому интерфейс системы должен быть интуитивным и предоставлять полный контроль над процессом планирования.